% Einheit 1 von 5 – Brüche addieren und subtrahieren (bank/bruchrechnung/e1.jsonl, zusammenbau v0.1)
%% TODO Zweigzeile Teil 1: was hier gelernt wird (Satz in Schülersprache aus dem Einheitstitel) – Einheit 1
\einheitenkopf[e1]{Einheit 1 von 5 · Brüche addieren und subtrahieren}
\zweigzeile{ab Kl. 5, je nach Buch bis Kl. 6 (am Gymnasium ab Kl. 6) · keine P10-Aufgabe}
\setcounter{aufgabe}{8}

%% TODO Titel: Ich-kann-Satz fehlt in der Bank (Platzhalter: Kettenname) – Nr. 9
%% TODO Anweisung: ein Satz über der ersten Teilaufgabe fehlt in der Bank – Nr. 9
\begin{aufgabe}{Gleicher Nenner – ja oder nein?}
\begin{teile}
\teil $\frac{4}{11} + \frac{5}{11}$ – sofort rechnen? \janein
\end{teile}
\end{aufgabe}

%% TODO Titel: Ich-kann-Satz fehlt in der Bank (Platzhalter: Kettenname) – Nr. 10
%% TODO Anweisung: ein Satz über der ersten Teilaufgabe fehlt in der Bank – Nr. 10
\begin{aufgabe}{Womit erweitere ich?}
\begin{teile}
\teil $\frac{1}{4}$ in Zwanzigstel – womit erweitern? mit \leerfeld
\end{teile}
\end{aufgabe}

%% TODO Titel: Ich-kann-Satz fehlt in der Bank (Platzhalter: Kettenname) – Nr. 11
%% TODO Anweisung: ein Satz über der ersten Teilaufgabe fehlt in der Bank – Nr. 11
\begin{aufgabe}{Addieren/Subtrahieren}
%% TODO \rechenplatz{n}: Schrittzahl des Grundfalls fehlt in der Bank (nur bei fester Schreibform, 2.3 b)
\begin{teile}
\teil Färbe $\frac{1}{6}$, dann weitere $\frac{4}{6}$ – wie viel ist gefärbt?

\bruchrechteck{0}{6}
\leerfeld
\teil $\frac{2}{11} + \frac{5}{11}$
\teil $\frac{5}{14} + \frac{3}{14}$
\teil $\frac{5}{7} + \frac{4}{7}$
\teil $\frac{1}{4} + \frac{3}{8}$
\teil $\frac{2}{3} + \frac{1}{4}$
\teil $2\frac{1}{5} + 1\frac{2}{5}$
\teil Ein Glücksrad hat 10 gleich große Felder: 5 rot, 3 blau, 2 gelb. Es wird zweimal gedreht. Die Pfade „zweimal rot“, „zweimal blau“ und „zweimal gelb“ haben die Wahrscheinlichkeiten $\frac{1}{4}$, $\frac{9}{100}$ und $\frac{4}{100}$. Wie groß ist die Wahrscheinlichkeit für zweimal dieselbe Farbe? Gib sie gekürzt an. \hfill (P10 2014 OS)
\teil Zwei Kreisel werden gedreht. Die Wahrscheinlichkeit, dass beide Rot zeigen, ist $\frac{4}{30}$. Wie groß ist die Wahrscheinlichkeit, dass nicht beide Rot zeigen? Gib sie gekürzt an. \hfill (P10 2026 FOR)
\end{teile}
\end{aufgabe}

%% TODO Titel: Ich-kann-Satz fehlt in der Bank (Platzhalter: Kettenname) – Nr. 12
%% TODO Anweisung: ein Satz über der ersten Teilaufgabe fehlt in der Bank – Nr. 12
\begin{aufgabe}{ganze Zahl plus Bruch}
\begin{teile}
\teil $3 + \frac{2}{5}$
\end{teile}
\end{aufgabe}

%% TODO Titel: Ich-kann-Satz fehlt in der Bank (Platzhalter: Kettenname) – Nr. 13
%% TODO Anweisung: ein Satz über der ersten Teilaufgabe fehlt in der Bank – Nr. 13
\begin{aufgabe}{Addieren/Subtrahieren – Fehler finden, Begründen, Anwendung, Darstellungswechsel}
\begin{teile}
\teil Jonas rechnet: $\frac{1}{5} + \frac{2}{5} = \frac{3}{10}$. Finde den Fehler.
\teil Warum darf man bei $\frac{1}{3} + \frac{1}{5}$ nicht einfach die Zähler addieren?
\teil Lina übt $\frac{3}{4}$ Stunde Klavier und danach $\frac{1}{2}$ Stunde Gitarre. Wie lange übt sie insgesamt?
\teil Im Streifen waren $\frac{7}{10}$ gefärbt. Jetzt ist nur noch so viel gefärbt wie im Bild. Welche Minusaufgabe passt, wie heißt das Ergebnis?

\bruchrechteck{4}{10}
\leerfeld
\end{teile}
\end{aufgabe}
